% =========================================================================
% SECCIÓN 6: CONCLUSIONES Y HOJA DE RUTA
% =========================================================================

\section{Conclusiones y Hoja de Ruta de la Investigación}
\label{sec:conclusiones_roadmap}

\subsection{Respuestas a las Preguntas de Investigación}
\begin{enumerate}[(a)]
	\item \textbf{Respuesta a RQ1 (SLMs Locales):} [Los modelos de la familia Qwen2.5 (1.5B--3B) cuantizados a 4 bits (GGUF) demuestran el mejor balance entre consumo de VRAM ($\le$3.5\,GB) y coherencia semántica en inferencia local offline.]
	\item \textbf{Respuesta a RQ2 (RAG Morfológico):} [La fragmentación guiada por unidades textuales con solapamiento morfológico reduce drásticamente el error de recuperación frente a chunking por longitud fija en Quechua Collao.]
	\item \textbf{Respuesta a RQ3 (Adaptación Eficiente):} [La adaptación vía QLoRA restringida a capas de atención específicas minimiza la degradación de capacidades lógicas del modelo base al tiempo que incorpora vocabulario pedagógico regional.]
	\item \textbf{Respuesta a RQ4 (Evaluación Integral):} [La métrica chrF++ correlaciona de manera superior a BLEU con el juicio de docentes bilingües debido a la naturaleza aglutinante del quechua.]
\end{enumerate}

\subsection{Aportes Teóricos y Prácticos al Plan de Tesis}
\noindent [Sustento teórico de la viabilidad de la propuesta y justificación técnica del stack tecnológico seleccionado.]

\subsection{Amenazas a la Validez y Limitaciones}
\noindent [Análisis de posibles sesgos en la búsqueda sistemática, dispersión de publicaciones no indizadas en lenguas originarias y limitaciones del hardware de experimentación local.]

\subsection{Hoja de Ruta (Roadmap) para el Desarrollo del Proyecto}
\begin{enumerate}[1.]
	\item \textbf{Fase 1: Curaduría y Normalización de Corpus.} [Recopilación de materiales MINEDU y textos Quechua Collao validados.]
	\item \textbf{Fase 2: Arquitectura RAG e Indexación Vectorial.} [Implementación de chunking morfológico y almacenamiento en base embebida.]
	\item \textbf{Fase 3: Optimización e Inferencia de SLM.} [Cuantización GGUF del modelo Qwen2.5-3B y pruebas de latencia en GTX 1050 Ti.]
	\item \textbf{Fase 4: Adaptación QLoRA y Alineación Pedagógica.} [Afinamiento de bajo rango para andamiaje educativo bilingüe.]
	\item \textbf{Fase 5: Evaluación Empírica y Validación con Expertos.} [Ejecución del protocolo de pruebas chrF++, latencia y rúbrica EIB.]
\end{enumerate}

\subsection{Líneas Futuras de Investigación}
\noindent [Proyecciones para el soporte multimodal (voz Quechua con Whisper-fine-tuned) y extensión a otras familias lingüísticas (Aymara, Asháninka).]

\subsubsection*{Agradecimientos}
\noindent [Agradecimiento institucional a la Universidad Nacional de San Antonio Abad del Cusco (UNSAAC), a la cátedra de Seminario de Investigación II bajo la dirección del Msc. Raúl Huillca Huallparimachi, y a los investigadores en PLN para lenguas originarias.]
